\documentclass[11pt]{article}
\usepackage[a4paper,margin=24mm]{geometry}
\usepackage{amsmath,amssymb,amsthm,microtype}
\usepackage[hidelinks]{hyperref}
\newcommand{\E}{\mathbb E}\newcommand{\Pp}{\mathbb P}
\newtheorem{theorem}{Theorem}\newtheorem{lemma}[theorem]{Lemma}
\title{The left probability range from five neighboring counts}
\author{}\date{28 October 2026}
\begin{document}\maketitle

Five neighboring counts retain enough of the crossing loss to pay for the
finite corrections. Centered replacements compare their smooth kernels
without an inward-displacement error. On the left probability range this
gives a complete signed exclusion at every index at least 700, including
the noise between the previous small-noise theorem and the geometric ceiling.

Write $\phi_0=(2\pi)^{-1/2}$ and $C_E=(3+\sqrt{10})\phi_0/6$.
\begin{theorem}\label{thm:left}
Let $n\ge700$ be an integer, $p\in[1/3,2/5]$, $q=1-p$, and
$Y=B+U$, where $B$ is Bernoulli of parameter $p$ and $\E(U\mid B)=0$.
Write $U_j$ for the conditional noise given $B=j$. Suppose
\[
 -7/10\le U_0\le a_0,\quad -a_1\le U_1\le7/10,
 \quad a_0,a_1\ge0,\quad a_0+a_1\le17/250.
\]
Put $v=pq$, $s=\E U^2$, $\lambda=ns$, $X=(Y-p)/\sqrt{v+s}$ and
$\rho=\E|X|^3$. If $1/100\le\lambda\le8/25$ and $|H|\le46/25$, then
for independent copies $X_1,\ldots,X_n$,
\[
 \sqrt n\{\Pp(X_1+\cdots+X_n\le H)-\Phi(H/\sqrt n)\}
                       -C_E\rho<-\lambda/100.
\]
\end{theorem}

There is no contact, predecessor or geometric-budget premise in this
comparison theorem. Those premises remain in acquiring the displayed strip
and noise ceiling for a putative extremal witness. The conclusion is signed;
reflection requires the inputs for the reflected law.

\section{Retaining the exact crossing loss}
Set $t=np+H\sqrt{v+s}$, $k=\lfloor t\rfloor$, $u=t-k\in[0,1)$,
$\delta=k-np$, and $s_j=\E U_j^2$. Since $v\in[2/9,6/25]$,
\[
 v+s<1/4,\quad |t-np|<23/25,\quad
 -48/25<\delta<23/25.
\]
In particular all five indices $k-1,\ldots,k+3$ are interior.
Use independent conditional pairs to form
\[
 T_j=j+\sum_{i\le j}U_{1i}+\sum_{i>j}U_{0i},\quad
 V_i=\sum_{l<i}U_{1l}+\sum_{l>i}U_{0l},
\]
\[
 h(x)=\E(x-U_0)_+-\E(x-1-U_1)_+-\Pp(1+U_1\le x),
 \qquad g_i=\E h(t-i+1-V_i).
\]
The mesh is increasing; $T_0\le t<T_n$ follows from the inward caps.
The finite crossing identity of v13 gives $g_i\ge0$ and, with
$\pi_j=\Pp(\operatorname{Bin}(n,p)=j)$,
\begin{equation}\label{eq:cross}
 \Pp(X_1+\cdots+X_n\le H)
 =(1-u)F_B(k)+uF_B(k+1)+\mathcal C(t)-\sum_i\pi_i g_i,
\end{equation}
where
\[
 \mathcal C(t)=\sum_{j=1}^{n-1}(\pi_{j+1}-\pi_j)
       \{\E(t-T_j)_+-(t-j)_+\}.
\]
Each conditional sum has the distribution from the actual iid law.
The auxiliary coupling changes none of these conditional marginals.
Lower crossings remain closed and upper crossings strict, including atoms.

\section{Five smooth kernels give the reserve}
Let $B(x)=\sin^2(\pi x)$ on $[0,1]$ and zero elsewhere. Its derivative
is globally Lipschitz with constant $2\pi^2$ and is zero at both joins.
We first check the pointwise minorant
\begin{equation}\label{eq:minorant} h(x)\ge(2/5)B(x).
\end{equation}
Separated support gives $h\ge0$. If $0\le x<1-a_1$, Jensen gives
$h(x)\ge x$. The elementary bound $B(x)\le(5/2)x$ follows from
$\pi^2<10$ for $x\le1/4$, the tangent at $1/4$ on $[1/4,2/5]$, and
$B\le1$ for $x\ge2/5$. On the middle interval $B$ is concave and its
tangent is $1/2+\pi(x-1/4)$; the two endpoint comparisons use $\pi<10/3$.

For $1-a_1\le x\le1$, put $D=1+U_1$ and $L=\E(D-x)_+$.
Put--call parity gives $h=\Pp(D>x)-L$, and hence
\[
 h(x)\ge\frac{x-7/10}{17/10-x}L
       \ge\frac{29}{96}(1-x),
\]
since $x\ge233/250$ and $L\ge1-x$. On this same interval,
\[
 (2/5)B(x)\le(2/5)(22/7)^2(17/250)(1-x)
                                    <(29/96)(1-x).
\]
At $x=1$ the minorant is zero. The strict survival probability in this
calculation retains all atoms of $D$.

Freeze $V=V_{k+1}$ and compare the five omitted-coordinate laws
$V^{(j)}\stackrel{\rm law}=V_{k+1+j}$, $-2\le j\le2$.
If $F'$ is Lipschitz with constant $L$ and $Z$ is centered, Taylor's
integral remainder gives
$|\E F(r+Z)-F(r)|\le L\E Z^2/2$. Thus one centered replacement
costs at most $\pi^2(s_0+s_1)$ for a bump with this derivative bound.
There are only four replacements after regrouping: for example, with
$B_j(x)=B(u-j-x)$, the two positive-index differences equal
\begin{align*}
 &\E(B_1+B_2)(V^{(1)})-\E(B_1+B_2)(V)\\
 &\hspace{15mm}+\E B_2(V^{(2)})-\E B_2(V^{(1)}).
\end{align*}
The open supports of $B_1,B_2$ are disjoint, and their derivatives
vanish at the common boundary. Their sum still has derivative Lipschitz
constant $2\pi^2$. The negative indices give two analogous replacements.
Since $s_0+s_1\le3\lambda/n$, their total cost is at most
$12\pi^2\lambda/n$. No independence of sums at different counts is assumed.

The infinite translates of $B$ sum to a periodic sine square. Every
omitted nonzero term with $|j|\ge3$ requires $|V|>2$, so
\[
 \sum_{j=-2}^2\E B(u-j-V)
 \ge\E\sin^2(\pi(u-V))-\E V^4/16.
\]
The two boundary bumps vanish at integer joins, so there is no missing
atom contribution.

For $|x|\le96/125$, one has
\begin{equation}\label{eq:cos}1-\cos(2\pi x)\ge(7/5)x^2.
\end{equation}
On $[0,3/4]$, the sine chord gives
$\sin(\pi x)\ge(2\sqrt2/3)x>(14/15)x$.
On $[3/4,96/125]$, its derivative bound gives
$\sin(\pi x)>7/10-(22/7)(9/500)=563/875$.
Squaring these positive bounds proves \eqref{eq:cos} on both closed
intervals, then evenness completes the proof.
Conditional support diameters are at most $7/10+17/250=96/125$.
For either conditional noise $Z$ and an independent copy $Z'$, this yields
\[
 1-|\E e^{2\pi iZ}|^2\ge(14/5)\operatorname{Var}Z,
 \qquad |\E e^{2\pi iZ}|\le e^{-(7/5)\operatorname{Var}Z}.
\]
Only the independent frozen summands are multiplied. Their variance is
\[
 \tau=\E V^2=(nq-\delta-1)s_0+(np+\delta)s_1,
 \qquad(1-6/n)\lambda\le\tau\le(1+3/n)\lambda.
\]
These are the physical count coefficients; they follow from the displayed
bounds for $\delta$, $p\ge1/3$ and $q\ge3/5$.
The fourth moment is at most $(49/100)\tau+3\tau^2$. Therefore
\begin{equation}\label{eq:gain}
 \sum_{i=k-1}^{k+3}g_i\ge\frac25\left\{
 \frac{1-e^{-(7/5)(1-6/n)\lambda}}2
 -\frac{(49/100)(1+3/n)\lambda+3(1+3/n)^2\lambda^2}{16}
 -\frac{12\pi^2\lambda}{n}\right\}.
\end{equation}
This is the left-strip adaptation of the four-replacement proof in
\path{loss-direct.tex}; the changes of support width, probabilities and
variance constants have all been derived here.

\section{Binomial weights and finite corrections}
For $|i-k-1|\le2$, the phase bounds give $|i-np|<4$.
The segment from $i/n$ to $p$ has variance product greater than $11/50$.
Finite Stirling with remainder $0<r_j<1/(12j)$ and the entropy Taylor
bound imply, uniformly for $n\ge700$,
\[
 \sqrt n\pi_i\ge2\phi_0
 \exp\left\{-\frac{8+1/12}{700(11/50)}\right\}
 >2\frac{19943}{50000}
       \left(1-\frac{8+1/12}{700(11/50)}\right)>\frac34.
\]
Thus \eqref{eq:gain} can be multiplied by $3/4$. The other count losses
are nonnegative and need no tail estimate.

Here is the summable curvature bound from v14, in its full form.
If $c=(np-q-t)_+\le1$, $T$ bounds the variances of the increasing-mass
counts, $M$ bounds their fourth moments, and
$\max_r\Pp(\operatorname{Bin}(n+1,p)=r)\le\beta/\sqrt{(n+1)v}$, then
\begin{equation}\label{eq:curvature}
 \sqrt n\mathcal C(t)\le\frac{\beta}{nv^{3/2}}
 \left\{T/4+c\sqrt T+\frac{27M}{256}(5/3+5c/4)\right\}.
\end{equation}
Indeed the adjacent binomial difference has numerator $np-q-j$.
For $a=t-j$, the centered stop-loss bend is bounded by $\sqrt T/2$,
by $T/(4a)$ when $a>0$, and by $27M/(256|a|^3)$ away from zero.
At most two nearby points are treated directly; the remaining full sums
use $\sum i^{-2}<5/3$ and $\sum i^{-3}<5/4$.
The increasing-mass indices satisfy $j<np$, so we may take
\[
 T=5\lambda/3,\quad M=(49/100)T+3T^2,\quad \beta=41/100.
\]
For completeness the last bound follows by finite Stirling at a mode:
for $N=n+1\ge43$, its variance product exceeds $39v/40$, whence the
mass is smaller than $(2/5)\sqrt{40/39}(516/515)/\sqrt{Nv}
<(41/100)/\sqrt{Nv}$. Also $v^{3/2}>13/125$, and
\[
                 \sqrt T\le(250/39)\lambda+13/200
\]
follows by squaring $\sqrt T-13/100$.

Schulz's strict finite-$n$ Bernoulli theorem gives at both integer anchors
$\sqrt n\{F_B(j)-\Phi((j-np)/\sqrt{nv})\}<C_E\rho_B$, where
$\rho_B=(1-2v)/\sqrt v$. Only this all-$p$ theorem is needed, not its
sharper restricted-probability version. Linear interpolation preserves
the strict inequality. Put $w=t-np$. Its normal interpolation error is at most
\[
 \frac{\phi_0}{2nv^{3/2}}u(1-u)(u-w)_+.
\]
This follows by bounding the positive part of the normal CDF's second
derivative over the whole unit interval.
For $H\le0$, $c<8/25$, the variance change is favorable and the remaining
normal error is at most
\[
                  \frac{2/5}{nv^{3/2}}(23/200+2/27).
\]
For $H\ge0$, $c=0$, and interpolation plus variance repair is at most
\[
                  \frac{2/5}{nv^{3/2}}(2/27+(23/50)\lambda).
\]
These use $|w|<23/25$, $u(1-u)\le1/4$, $u^2(1-u)\le4/27$ and
$1-(1+x)^{-1/2}\le x/2$.

The actual cubic correction is favorable:
\[
 (v+s)^{3/2}(\rho-\rho_B)
 \ge s\{1-5/6-17/250-(15/16)s\}>0,
\]
since $s\le(8/25)/700$. To obtain this bound, expand the marginal absolute
third moment and retain its signs, use
$\E U_0^3\le a_0s_0$, $\E U_1^3\ge-a_1s_1$, and
$(1+x)^{3/2}\le1+3x/2+3x^2/8$. No favorable cubic sign is inferred
below $p=1/3$.

\section{Two concave quartics close the entire noise interval}
Set $N=700$, $a=1-6/N$, $b=1+3/N$, $z=(7/5)a$ and define
\[
 W(\lambda)=\frac3{10}\left\{
 \frac{z\lambda-z^2\lambda^2/2+z^3\lambda^3/6-z^4\lambda^4/24}{2}
 -\frac{(49/100)b\lambda+3b^2\lambda^2}{16}
 -\frac{12(22/7)^2\lambda}{N}\right\}.
\]
Taylor's theorem gives $1-e^{-x}\ge x-x^2/2+x^3/6-x^4/24$ for $x\ge0$.
Every term in \eqref{eq:gain} improves with $n$, so $W$ is a lower bound
for the normalized weighted loss at every $n\ge N$.
Write
\[
 S_c=T/4+c\{(250/39)\lambda+13/200\}
                  +(27M/256)(5/3+5c/4),
\]
\[
 E_-(\lambda)=\frac{(41/100)S_{8/25}+(2/5)(23/200+2/27)}{13/125},
\]
\[
 E_+(\lambda)=\frac{(41/100)S_0+(2/5)(2/27+(23/50)\lambda)}{13/125}.
\]
These are nonnegative quadratics. The preceding bounds reduce the theorem
to $P_\pm(\lambda)>0$ on $[1/100,8/25]$, where
\[
                  P_\pm=W-E_\pm/N-\lambda/100.
\]
The gain polynomial is strictly concave on the whole real line, since
\[
 W''(\lambda)=-\frac3{10}\left\{
 \frac{z^2}{4}\bigl((z\lambda-1)^2+1\bigr)+\frac{3b^2}{8}\right\}<0.
\]
Both error quadratics are convex. Hence $P_\pm''<0$, and their minima
on the closed noise interval occur at endpoints. The four exact endpoint
margins are positive. The smallest is
\[
 \frac{5442273179742877948871}{97819312500000000000000000}>0.
\]
The companion rational checker reconstructs all coefficients and guards;
it does not accept sampled positivity or a fitted grid. This proves
Theorem~\ref{thm:left} with all phases and noise endpoints included.

\section*{Scope and evidence}
The kernel and four-replacement method are in this round's
\path{loss-direct.tex}; the v13 crossing identity and v14 summable-curvature
proof retain their stated atom conventions. The external Bernoulli theorem
is Schulz, \emph{The optimal Berry--Esseen constant in the binomial case}
(2016), Theorem 1, printed page 1,
\href{https://d-nb.info/1197702695/34}{original thesis}.
Finite Stirling and the normal constants are explicitly enclosed in the
checker. No Fourier or certificate-generation campaign is used.

Combined with the retained LEFT-01 comparison below $\lambda=.01$, this
closes the acquired $p\in[1/3,.4]$ strip at all $n\ge700$ through its noise
ceiling. The retained geometry excludes $p<1/3$ already at $n\ge169$.
Acquisition, reflection, attainment, remaining finite indices and BASE41
remain separate obligations. No unconditional cutoff theorem is claimed.
\end{document}
